\begin{notechapter}
\chapter{A Neural Bigram Model and the Embedding Matrix}\label{ch:bigram-neural}
\chaptersubtitle{A hidden layer between the previous word and the next: lookup, bottleneck and sharing}

\begin{chapterintro}
The softmax regression of \cref{ch:bigram-models} has one free row of logits per context: it
cannot transfer anything it learns about the word ``cat'' to the word ``dog''. Inserting a
hidden layer of $D\ll V$ units between the one-hot input and the softmax changes this. The first
layer becomes a table of $D$-dimensional vectors, one per word, the \emph{embeddings}; the second
layer must predict the next word from these vectors only, so words that are followed by
similar words are pushed towards similar vectors. This chapter derives the gradients of
the model, corrects a numerical example, and makes precise what the bottleneck
constrains. You will learn
\begin{itemize}[nosep]
  \item the one-hidden-layer bigram model and its parameter count;
  \item why its first layer is an embedding lookup, as in \texttt{nn.Embedding};
  \item the complete backward pass, with every shape checked, on a worked example verified
        with PyTorch;
  \item the rank constraint imposed by the hidden layer, and what it buys and costs.
\end{itemize}
\prerequisites{\cref{ch:bigram-models}; the backpropagation algorithm (\cref{ch:backpropagation}).}
\end{chapterintro}

\section{The model}\label{bigram-neural:sec:model}

The input is the one-hot code $x = e_i\in\R^V$ of the previous word. With $W_1\in\R^{V\times D}$, $b_1\in\R^D$, $W_2\in\R^{D\times V}$
and $b_2\in\R^V$,
\[
  a = W_1^\top x + b_1,\qquad h = \tanh(a),\qquad z = W_2^\top h + b_2,\qquad p = \softmax(z),
\]
and the loss of the bigram $(i, j)$ is $\ell = -\log p_j$. As in \cref{ch:bigram-models} the weights multiply their
input transposed, $W^\top x$ (the row-vector convention \texttt{x @ W} of the original notes and of
the code below).

\begin{figure}[ht]
\centering
\begin{tikzpicture}[
  box/.style={draw, rounded corners=2pt, minimum height=8mm, inner sep=3pt, font=\small, align=center},
  ar/.style={-{Stealth[length=1.8mm]}, thick},
  lab/.style={font=\scriptsize, inner sep=1pt}]
  \node[box, fill=bookaccent!15] (x) at (0,0) {$x = e_i$\\[-1pt]{\scriptsize one-hot, $V$}};
  \node[box, fill=bookexample!15] (h) at (4.3,0) {$h = \tanh(W_1^\top x + b_1)$\\[-1pt]{\scriptsize $D$ units}};
  \node[box, fill=bookkey!12] (z) at (8.9,0) {$z = W_2^\top h + b_2$\\[-1pt]{\scriptsize logits, $V$}};
  \node[box, fill=gray!10] (p) at (12.0,0) {$p = \softmax(z)$};
  \draw[ar] (x) -- node[lab, above]{$W_1$: $V\times D$} (h);
  \draw[ar] (h) -- node[lab, above]{$W_2$: $D\times V$} (z);
  \draw[ar] (z) -- (p);
  \draw[decorate, decoration={brace, mirror, amplitude=4pt}] (-0.8,-0.75) -- node[lab, below=5pt]{row $i$ of $W_1$: the embedding of word $i$} (6.1,-0.75);
\end{tikzpicture}
\caption{The neural bigram model. The first layer reads one row of $W_1$; the second must predict the
next word from the $D$ numbers of $h$.}
\label{bigram-neural:fig:model}
\end{figure}

\paragraph{The first layer is a lookup.} As in \cref{bigram-models:sec:softmax}, $W_1^\top e_i$ is row $i$ of $W_1$, transposed.
The hidden pre-activation is $a = (W_1)_{i\cdot}^\top + b_1$: the first layer stores one vector of $D$ numbers per word,
the \emph{embedding} of the word, and the one-hot product only retrieves it. This is exactly
\texttt{nn.Embedding(V, D)} in PyTorch, which stores the $V\times D$ matrix and indexes it, so that
the $V\times D$ product with a one-hot vector is never computed. As before, $b_1$ is redundant with a one-hot
input ($W_1 + \mathbf 1b_1^\top$ absorbs it). The second-layer bias $b_2$ is not: it adds the same vector to the
logits of every context, which the rows of $W_2^\top h$ can do only if some combination of hidden units is
constant across words.

\paragraph{Parameters.} The model has $VD + D + DV + V$ parameters, against $V^2$ for the table of
\cref{ch:bigram-models}. For $V = 10^4$ and $D = 100$: about $2\times10^6$ instead of $10^8$.

\section{The backward pass}\label{bigram-neural:sec:backward}

The three local rules of \cref{backpropagation:sec:local} give, from the output down,
\begin{align*}
  \bar z &= p - e_j &&\text{(softmax and cross-entropy)}\\
  \nabla_{W_2}\ell &= h\,\bar z^\top\in\R^{D\times V},\qquad\nabla_{b_2}\ell = \bar z &&\text{(affine map }z = W_2^\top h + b_2\text{)}\\
  \bar h &= W_2\,\bar z\in\R^D &&\text{(same affine map, with respect to its input)}\\
  \bar a &= \bar h\odot(1 - h\odot h) &&\text{(}\tanh' = 1 - \tanh^2\text{)}\\
  \nabla_{W_1}\ell &= x\,\bar a^\top\in\R^{V\times D},\qquad\nabla_{b_1}\ell = \bar a &&\text{(affine map }a = W_1^\top x + b_1\text{)}.
\end{align*}
Because the layers use $W^\top$, the outer products are (input)$\times$(gradient at the output)$^\top$, the
transposes of those of \cref{backpropagation:thm:backprop}, and each gradient has the shape of its
matrix. With $x = e_i$, $\nabla_{W_1}\ell$ has the single non-zero row $i$, equal to $\bar a^\top$: one example updates only the
embedding of its own context word, but it updates \emph{all} of $W_2$, which is shared by all
contexts.

\begin{pitfall}[Shapes and the factor $1 - h^2$]
Two mistakes are common in hand-written backward passes for this model, and both were
in the original version of this chapter. The first is a shape error: writing $\nabla_{W_2}\ell = \bar z\,h^\top$ ($V\times D$)
for a matrix $W_2$ of shape $D\times V$; the update $W_2 - \eta\nabla_{W_2}\ell$ is then not even defined, and a hand calculation
silently mixes entries. The second is applying the derivative of $\tanh$ twice, once in $\bar a$ and
again in $\nabla_{W_1}\ell$. Checking every gradient's shape against its parameter, and comparing with
automatic differentiation, catches both.
\end{pitfall}

\section{A worked example}\label{bigram-neural:sec:example}

\begin{example}[One step, verified]\label{bigram-neural:ex:step}
Vocabulary (cat, dog, mouse), $V = 3$, $D = 2$; the bigram is (cat, dog), so $x = e_1$ and $j = 2$. The parameters are
\[
  W_1 = \begin{bmatrix}0.1 & 0.2\\ 0.3 & 0.4\\ 0.5 & 0.6\end{bmatrix},\quad b_1 = \begin{bmatrix}0.1\\ 0.2\end{bmatrix},\quad
  W_2 = \begin{bmatrix}0.7 & 0.8 & 0.9\\ 1.0 & 1.1 & 1.2\end{bmatrix},\quad b_2 = \begin{bmatrix}0.3\\ 0.4\\ 0.5\end{bmatrix}.
\]
All numbers are rounded to four digits.

\emph{Forward.} $a = (0.1, 0.2) + (0.1, 0.2) = (0.2, 0.4)$ (row ``cat'' of $W_1$ plus $b_1$), $h = \tanh(a) = (0.1974, 0.3799)$, and
\[
  z = W_2^\top h + b_2 = \begin{bmatrix}0.7\cdot0.1974 + 1.0\cdot0.3799 + 0.3\\ 0.8\cdot0.1974 + 1.1\cdot0.3799 + 0.4\\ 0.9\cdot0.1974 + 1.2\cdot0.3799 + 0.5\end{bmatrix}
  = \begin{bmatrix}0.8181\\ 0.9758\\ 1.1336\end{bmatrix} .
\]
Then $e^z = (2.2661, 2.6533, 3.1069)$, whose sum is $8.0263$, so $p = (0.2823, 0.3306, 0.3871)$ and $\ell = -\log0.3306 = 1.1069$.

\emph{Backward.} $\bar z = p - e_2 = (0.2823, -0.6694, 0.3871)$, and
\[
  \nabla_{W_2}\ell = h\,\bar z^\top = \begin{bmatrix}0.0557 & -0.1321 & 0.0764\\ 0.1073 & -0.2543 & 0.1471\end{bmatrix},\qquad \nabla_{b_2}\ell = \bar z .
\]
The gradient reaching the hidden layer is
$\bar h = W_2\bar z = (0.1976 - 0.5355 + 0.3484,\ 0.2823 - 0.7363 + 0.4645) = (0.0105,\ 0.0105)$. With $1 - h\odot h = (0.9610,\ 0.8556)$,
$\bar a = (0.0101,\ 0.0090)$, and
\[
  \nabla_{W_1}\ell = x\,\bar a^\top = \begin{bmatrix}0.0101 & 0.0090\\ 0 & 0\\ 0 & 0\end{bmatrix},\qquad\nabla_{b_1}\ell = \bar a .
\]
\emph{Update} with $\eta = 0.1$: only row ``cat'' of $W_1$ changes, to $(0.0990, 0.1991)$; $b_1$ becomes $(0.0990, 0.1991)$,
$b_2 = (0.2718, 0.4669, 0.4613)$ and
\[
  W_2 = \begin{bmatrix}0.6944 & 0.8132 & 0.8924\\ 0.9893 & 1.1254 & 1.1853\end{bmatrix}:
\]
the column of ``dog'' grows, the two others shrink. The hidden gradient $\bar h$ is small because the
three columns of $W_2$ are nearly parallel: the hidden layer has little influence on which word
wins, and almost all the learning happens in the output layer at this point.
\end{example}

The same computation in NumPy, compared with PyTorch's automatic differentiation:
\begin{lstlisting}[language=Python]
import numpy as np, torch
W1 = np.array([[0.1, 0.2], [0.3, 0.4], [0.5, 0.6]]); b1 = np.array([0.1, 0.2])
W2 = np.array([[0.7, 0.8, 0.9], [1.0, 1.1, 1.2]]); b2 = np.array([0.3, 0.4, 0.5])
i, j = 0, 1                                    # previous word 'cat', next word 'dog'
h = np.tanh(W1[i] + b1)                        # W1.T @ one_hot(i) is the row W1[i]
z = W2.T @ h + b2
p = np.exp(z - z.max()); p /= p.sum()
dz = p.copy(); dz[j] -= 1.0                    # dloss/dz = p - e_j
dW2, db2 = np.outer(h, dz), dz                 # shape (D, V), like W2
da = (W2 @ dz) * (1 - h**2)                    # tanh' applied once
dW1 = np.zeros_like(W1); dW1[i] = da; db1 = da # only row i of W1

t = {k: torch.tensor(v, requires_grad=True) for k, v in
     dict(W1=W1, b1=b1, W2=W2, b2=b2).items()}
zt = t["W2"].T @ torch.tanh(t["W1"][i] + t["b1"]) + t["b2"]
(-torch.log_softmax(zt, 0)[j]).backward()
assert np.allclose(t["W1"].grad, dW1) and np.allclose(t["W2"].grad, dW2)
\end{lstlisting}

\section{What the bottleneck does}\label{bigram-neural:sec:bottleneck}

Stack the hidden vectors of all the words as the rows of $H = \tanh(W_1 + \mathbf 1b_1^\top)\in\R^{V\times D}$. The matrix of logits of all
contexts, row $i$ for the previous word $i$, is
\[
  Z = HW_2 + \mathbf 1b_2^\top,\qquad\rank Z\leq D + 1 .
\]
The table model of \cref{ch:bigram-models} allows any $V\times V$ matrix of logits; the neural model only those
of rank at most $D + 1$. This constraint is the point of the model.
\begin{itemize}
  \item \emph{It forces sharing.} Two words whose rows of $\hat p(\cdot\mid i)$ are similar can only be fitted
        cheaply by giving them similar hidden vectors, since the same $W_2$ maps every $h$ to logits.
        The model then predicts similar continuations for them, including continuations
        that were observed after one and never after the other. This is how a neural
        language model assigns non-zero probability to unseen bigrams without smoothing
        (Bengio, Ducharme, Vincent and Jauvin, \emph{A Neural Probabilistic Language Model},
        2003).
  \item \emph{It needs data to act.} If $D + 1$ is large enough to fit the empirical table, the model
        can reproduce the counts and inherits their zeros, like the table model.
\end{itemize}

\begin{example}[The four-sentence corpus with $D = 2$]\label{bigram-neural:ex:corpus}
Train the model with \texttt{nn.Embedding(8, 2)}, $\tanh$ and \texttt{nn.Linear(2, 8)} on the corpus of
\cref{bigram-models:ex:corpus} (Adam, learning rate $0.05$, $3\,000$ steps; five random seeds). In every run the
training perplexity is $1.462$, that of the count model: the empirical table of this tiny corpus
is simple enough to be fitted with rank $3$. The hidden vectors show the sharing: ``sat'' and
``ran'', which are both always followed by $\langle/s\rangle$, end with almost the same $h$ in every run (for
instance $(0.99, 1.00)$ and $(1.00, 0.98)$). But the held-out sentence \emph{a dog ran} still has
perplexity between $40$ and $54$, with $p(\text{dog}\mid\text{a})\approx0.001$: four sentences do not contain enough evidence
that ``a'' behaves like ``the''. Sharing generalizes only what the data make cheaper to share.
\end{example}

\begin{intuition}
The table model learns $V$ unrelated distributions. The neural model learns a map from a
$D$-dimensional space of ``contexts'' to distributions, and places each word in that space.
Words end up close when they predict similar continuations; their positions are the
embeddings that \cref{ch:word-embeddings} studies for their own sake.
\end{intuition}

\paragraph{Cost.} Each example costs $O(D)$ for the lookup but $O(DV)$ for the logits and the softmax: the
output layer dominates as soon as the vocabulary is large. The next chapters are about this
cost: hierarchical softmax (\cref{ch:hierarchical-softmax}) and negative sampling
(\cref{ch:nce-negative-sampling}) avoid computing the $V$ logits.

\begin{chaptersummary}
\begin{itemize}
  \item The neural bigram model computes $h = \tanh(W_1^\top e_i + b_1)$, $z = W_2^\top h + b_2$, $p = \softmax(z)$, with $2VD + D + V$
        parameters instead of $V^2$.
  \item Its first layer is a table of word vectors read by index: an embedding layer. With a
        one-hot input, $b_1$ is redundant.
  \item Backward pass: $\bar z = p - e_j$, $\nabla_{W_2}\ell = h\bar z^\top$, $\bar a = (W_2\bar z)\odot(1 - h^2)$, $\nabla_{W_1}\ell = x\bar a^\top$ (only the row of the
        context word); every gradient has the shape of its parameter.
  \item The logits of all contexts form a matrix of rank at most $D + 1$; the constraint forces
        words with similar continuations to share their representation, which is how neural
        language models generalize to unseen bigrams when there is enough data.
\end{itemize}
\end{chaptersummary}

\section*{Exercises}
\addcontentsline{toc}{section}{Exercises}

\begin{exercise}\label{bigram-neural:exr:1}
Check that every gradient of \cref{bigram-neural:ex:step} has the shape of its parameter, and
recompute the entry $(2, 3)$ of $\nabla_{W_2}\ell$.
\solutionlink{bigram-neural:exr:1}
\end{exercise}

\begin{exercise}\label{bigram-neural:exr:2}
The original code of this chapter computed \texttt{dh = W2 @ dz * (1 - h**2)} and then
\texttt{dW1 = np.outer(x, dh * (1 - h**2))}. What does it compute instead of $\nabla_{W_1}\ell$? Give its non-zero row in
the example.
\solutionlink{bigram-neural:exr:2}
\end{exercise}

\begin{exercise}\label{bigram-neural:exr:3}
Show that $\rank(HW_2 + \mathbf 1b_2^\top)\leq D + 1$. With $D = 1$ and no biases, can the model represent a bigram table
in which word $1$ is always followed by word $2$ and word $2$ always by word $1$, while word $3$ is
followed by each word with probability $1/3$?
\solutionlink{bigram-neural:exr:3}
\end{exercise}

\begin{exercise}\label{bigram-neural:exr:4}
Count the parameters for $V = 50\,000$ and $D = 300$, and the multiply-adds per example of the lookup and
of the output layer.
\solutionlink{bigram-neural:exr:4}
\end{exercise}
\end{notechapter}
